\section{Introduction}
\label{sec:introduction}

Generative-AI products face growing litigation over training data, output reproduction, model behavior, and downstream misuse \citep{henderson2023foundation,levendowski2018copyright,grimmelmann2015copyright}.
High-profile suits such as \textit{Getty Images v.\ Stability AI}~\citep{getty_stability_2023}, \textit{New York Times v.\ Microsoft}~\citep{nyt_openai_2023}, and \textit{Kadrey v.\ Meta}~\citep{kadrey_meta_2023} illustrate the range of theories from DMCA §1202 information removal to direct infringement and vicarious liability for outputs.
For the engineering organizations that build these systems, the operational question is not merely whether a lawsuit exists, but whether a theory of liability maps onto a specific product attribute, whether the supporting authority remained valid at the relevant date, and which mitigation should be reviewed by counsel.
Legal-operations teams track dockets and news, yet the bridge from a litigation event to a product-specific engineering action is usually manual.
Keyword alerts surface noisy filings, docket trackers rarely map cases to product features, and conventional retrieval-augmented generation (RAG) \citep{lewis2020retrieval} treats semantically similar legal documents as interchangeable even when later decisions have narrowed or superseded them.

We formalize this need as Product--Litigation Risk Entailment (\plre): deciding whether product attributes entail exposure to a legal-risk theory under time-valid precedent.
\plre{} differs from standard legal textual entailment \citep{rabelo2022coliee} in three respects: the premise is an engineering description rather than legal text, the relation is directional (product attribute $\to$ risk theory), and the support set is conditioned on which authority was operative at the query date.
Unlike commercial citators such as Shepard's and KeyCite, we construct a public-corpus, automatically extracted, time-indexed treatment layer and connect it directly to product-facing entailment.

We present \sysname{} (Validity-Interval Gating in Litigation), a continuous sensing system built around this task.
Its core component is a supersession-aware treatment layer that infers when each authority ceases to be operative from treatment edges in later opinions.
A simple retrieval-time gate applies these intervals before generation, rather than flagging stale citations after an LLM has already produced a report.
Around this layer, a goal-conditioned agentic loop exposes unmet evidence requirements and stops when a coverage threshold is met.
Although the treatment-layer construction is in principle applicable to any time-stamped corpus in which later documents can supersede earlier ones, we evaluate it only for U.S.\ case law within \sysname{}; external-stack transfer remains a pilot and is not claimed.

Our contributions are:
\begin{itemize}[leftmargin=*]
\item a supersession-aware treatment layer that extracts treatment edges from public opinions, validates them against gold spans and a commercial citator, and propagates them into validity intervals that gate retrieval at query time (\S\ref{sec:temporal});
\item \plre{}, an initial product-conditioned legal-risk entailment benchmark with non-VIGIL baselines and ablations that separate temporal validity from flat similarity (\S\ref{sec:plre}, \S\ref{sec:experiments});
\item a goal-conditioned agentic loop with a coverage-based stopping rule; a conformal omission-risk bound holds under exchangeability and is reported as an auxiliary control (it is violated under temporal shift; \S\ref{sec:goal}, \S\ref{sec:experiments});
\item VIGIL-Dataset with per-record provenance and a human-adjudicated subset (\S\ref{sec:dataset}), plus a public SCOTUS overruling pair resource (see supplementary material).
\end{itemize}
